\chapter{CONCLUSION AND FUTURE WORK}
\label{ch:conclusion_future_work}

This chapter concludes the study on \texttt{kfc-procedure: A Unified Framework for Clusterwise Learning and COBRA-Based Ensemble Aggregation}. It summarizes the implemented framework, identifies the main contributions supported by the repository, discusses the limitations observed from the implementation and experiment notebooks, and proposes future work grounded in the current system.

\section{Conclusion}
\label{sec:conclusion}

This thesis presented the design and implementation of the \texttt{kfc\_procedure} framework, a Python-based machine learning library that integrates clusterwise learning with COBRA-based ensemble aggregation. The framework was developed as a modular research package that allows clustering, local supervised learning, and prediction aggregation to be configured and combined within a unified workflow.

The first major component of the framework is the KFCProcedure clusterwise learning pipeline. This pipeline is organized into three main stages: the K-step, the F-step, and the C-step. The K-step performs divergence-based clustering using Bregman K-Means. The repository implements several Bregman divergence choices, including squared Euclidean, generalized KL divergence, Itakura--Saito divergence, and logistic loss divergence. These divergences allow the same dataset to be partitioned under different geometric assumptions.

The F-step performs clusterwise supervised learning. Instead of fitting only one global estimator, the framework trains local models for each divergence-specific cluster. During prediction, each sample is assigned to its corresponding cluster under each divergence and routed to the appropriate local model. This produces a prediction matrix in which each column represents one divergence-level predictive view.

The C-step performs ensemble aggregation. It receives the prediction matrix produced by the F-step and combines the divergence-level predictions into a final output. The implemented C-step combiners include simple averaging, weighted mean aggregation, stacking-based regression and classification, majority voting, and COBRA-based aggregation strategies. This design allows the final aggregation mechanism to be changed without modifying the clustering or local learning stages.

The second major component of the repository is the COBRA aggregation subsystem. The implementation includes \texttt{GradientCOBRA} for regression aggregation in prediction space, \texttt{MixCOBRARegressor} for mixed input-space and prediction-space aggregation, and \texttt{CombinedClassifier} for classification aggregation through weighted voting. The COBRA subsystem also provides reusable components for distance computation, kernel weighting, loss functions, optimization, cross-validation, normalization, adapters, splitters, and aggregators.

From a software design perspective, the repository demonstrates modularity, separation of concerns, registry-based component management, and scikit-learn-style estimator design. The framework separates clustering, local model training, aggregation, optimization, and utility functions into distinct modules. It also uses factory registries to support extensibility for divergences, local models, combiners, kernels, distances, losses, optimizers, splitters, adapters, normalizers, and aggregators.

The experimental scope is defined by the notebooks under \texttt{notebooks/experiments/}. The classification experiments include Heart Disease, Diabetes, and Loan Prediction datasets. The regression experiments include Abalone, Housing Price, and Wind Turbine datasets. The notebooks evaluate baseline machine learning models, COBRA-based models, and KFC configurations where implemented. The results show that the framework components can be executed on multiple supervised learning tasks. However, the provided experimental outputs do not support a general claim that KFCProcedure consistently outperforms all baseline models. Therefore, this conclusion does not claim overall performance superiority.

\section{Key Contributions}
\label{sec:key_contributions}

The main contributions of this work are limited to what is directly supported by the repository, source code, and notebooks.

\begin{enumerate}
\item \textbf{Implementation of a unified KFC pipeline.}
The repository implements the \texttt{KFCProcedure} framework as an end-to-end pipeline composed of K-step clustering, F-step local model training, and C-step prediction aggregation.

```
\item \textbf{Divergence-based clusterwise learning.}  
The K-step implements Bregman K-Means and supports multiple divergence functions. This allows the framework to generate different cluster partitions of the same input space according to different divergence assumptions.

\item \textbf{Cluster-local supervised modeling.}  
The F-step trains local estimators for each divergence-specific cluster. This provides an implementation of clusterwise learning where supervised models are specialized to local regions of the data.

\item \textbf{Flexible aggregation through the C-step.}  
The C-step supports several aggregation strategies, including mean aggregation, weighted mean aggregation, stacking, majority voting, and COBRA-based combiners.

\item \textbf{Implementation of COBRA-based aggregation mechanisms.}  
The repository includes \texttt{GradientCOBRA}, \texttt{MixCOBRARegressor}, and \texttt{CombinedClassifier}. These components implement prediction-space or mixed-space aggregation using distances, kernels, optimization, and weighted aggregation.

\item \textbf{Modular component architecture.}  
The codebase separates major responsibilities into packages for clustering, local models, workflow steps, combiners, COBRA core components, and utilities. This improves traceability between the methodology and implementation.

\item \textbf{Factory and registry-based extensibility.}  
The repository implements factory systems for registering and resolving framework components. This supports extensibility without requiring changes to the main orchestration logic.

\item \textbf{Scikit-learn-style interface design.}  
The main predictive classes follow the common \texttt{fit()} and \texttt{predict()} workflow used by scikit-learn estimators. Some classes also inherit from scikit-learn base classes and mixins, which supports familiar usage patterns for machine learning experiments.

\item \textbf{Experiment notebooks for classification and regression tasks.}  
The repository includes notebooks demonstrating classification and regression experiments using baseline models, COBRA models, and KFC configurations where available.
```

\end{enumerate}

These contributions show that the project implements a reusable research framework for studying the interaction between clusterwise learning and ensemble aggregation. The contributions are primarily architectural and methodological rather than claims of universal predictive superiority.

\section{Limitations}
\label{sec:limitations}

Although the repository implements the main components of the proposed framework, several limitations are observable from the codebase and notebooks.

\subsection{Experimental Coverage Limitations}

The notebooks provide model comparison experiments, but they do not contain a complete formal ablation study. In particular, no systematic experiment is provided to isolate the effect of each divergence function, the number of clusters, the local estimator choice, or the aggregation method under controlled conditions.

The classification notebooks provide completed repeated experiments for the \texttt{KFCClassifier}. However, completed KFC regression results are not fully reported in the provided notebooks. The Abalone notebook contains an attempted KFC regression experiment, but the run is interrupted during the \texttt{kfc-gradientcobra} configuration. The Housing Price and Wind Turbine notebooks do not report completed KFC regression experiments.

Therefore, the experimental evidence supports framework demonstration and model comparison, but it does not fully establish the conditions under which KFCProcedure is expected to outperform baseline models.

\subsection{Testing Limitations}

The automated tests in the repository focus mainly on the COBRA core components, including splitters, estimators, normalizers, distances, kernel adapters, kernels, losses, cross-validation, optimizers, and aggregators. No separate automated test suite is observed for the full KFC subsystem, including \texttt{BregmanKMeans}, \texttt{KStep}, \texttt{FStep}, \texttt{CStep}, and the complete \texttt{KFCProcedure} training and prediction workflow.

As a result, the current test suite provides useful validation for the COBRA core layer, but full system-level validation of the KFC pipeline remains incomplete.

\subsection{Probability Prediction Limitation}

The repository defines a \texttt{predict\_proba()} method in the high-level KFC interface. However, the current KFC prediction workflow calls \texttt{FStep.predict\_proba()}, and this method is not implemented in the observed \texttt{FStep} class. Therefore, full probability prediction through the complete KFC pipeline should not be considered a completed feature in the current implementation.

\subsection{Scalability and Computational Limitations}

The framework includes some implementation choices that support larger computations, such as block-wise distance computation in Bregman clustering and optional acceleration paths in the COBRA classification subsystem. However, the repository does not provide a formal scalability benchmark or distributed training implementation.

The COBRA-based methods rely on distance matrices and cross-validation-based parameter selection, which may become computationally expensive for larger datasets. This limitation is also suggested by the interrupted KFC regression experiment in the Abalone notebook, although no formal timing or memory analysis is reported.

\subsection{Preprocessing and Workflow Limitations}

The main KFC orchestration layer expects the user or notebook to provide appropriately preprocessed data. General preprocessing steps such as missing-value imputation, categorical encoding, and feature scaling are implemented in the notebooks rather than inside the KFCProcedure pipeline itself.

This design keeps the framework flexible, but it also means that users must ensure that input data satisfy the domain requirements of the selected divergences. For example, generalized KL and Itakura--Saito divergences require positive-valued inputs, while the logistic divergence requires values in the interval $(0,1)$.

\subsection{Documentation and CI/CD Limitations}

The repository contains README documentation, notebooks, package metadata, and a tag-triggered GitHub Actions workflow for testing, building, and publishing. However, no separate documentation-generation system, such as Sphinx or MkDocs, is observed. The CI/CD workflow is also not configured as a broad every-commit validation pipeline with linting, static analysis, and notebook execution.

Thus, documentation and automation exist, but they remain limited compared with a fully mature production package workflow.

\section{Future Work}
\label{sec:future_work}

Several future improvements can be developed based on the current repository structure and observed limitations.

\subsection{Extension of the K-Step}

The K-step can be extended by adding more divergence families and clustering strategies. The existing Bregman divergence factory already provides a natural extension point for new divergence implementations. Future work may include additional Bregman divergences, adaptive divergence selection, or dataset-specific divergence configuration.

Another useful direction is systematic cluster-number selection. The current notebooks use fixed values of \texttt{n\_clusters}. Future experiments could evaluate multiple cluster numbers and select the best configuration through validation metrics or internal clustering criteria.

\subsection{Improvement of the F-Step}

The F-step can be improved by expanding the set of local estimators and improving how local models are selected. The current factory system already supports scikit-learn-style estimators, and this can be extended to include additional regression and classification models.

Future work may also explore automatic local model selection, where different clusters are allowed to use different estimator types depending on local data characteristics. This would further strengthen the clusterwise learning objective of the framework.

\subsection{Enhancement of the C-Step}

The C-step can be extended with additional aggregation methods. Current combiners include simple aggregation, stacking, and COBRA-based aggregation. Future work may explore more robust aggregation strategies, adaptive weighting, calibration-aware aggregation, and uncertainty-aware aggregation.

The COBRA-based combiners may also be improved by expanding kernel choices, improving bandwidth selection, and adding more efficient nearest-neighbor search methods for large prediction matrices.

\subsection{Hyperparameter Optimization}

The current notebooks configure divergences, cluster numbers, local models, and combiners manually. Future work should introduce a more systematic hyperparameter optimization workflow. This may include grid search, randomized search, Bayesian optimization, or nested cross-validation for selecting:

\begin{itemize}
\item divergence sets;
\item number of clusters;
\item local model types;
\item combiner types;
\item COBRA kernel functions;
\item COBRA distance functions;
\item optimizer and loss configurations.
\end{itemize}

Such a workflow would make the framework easier to compare under controlled and reproducible experimental conditions.

\subsection{Scalability Improvements}

Future work should address computational efficiency for large datasets. Possible improvements include:

\begin{itemize}
\item parallel training of local cluster models;
\item parallel computation of divergence-specific cluster assignments;
\item approximate nearest-neighbor search for COBRA aggregation;
\item memory-efficient distance matrix computation;
\item optional GPU acceleration for distance and kernel operations;
\item distributed execution for large-scale experiments.
\end{itemize}

These improvements would be especially useful for COBRA-based aggregation, where distance computation and parameter optimization can become expensive.

\subsection{Expanded Testing and Benchmarking}

The current automated tests focus on COBRA core components. Future work should add unit and integration tests for:

\begin{itemize}
\item Bregman divergence implementations;
\item \texttt{BregmanKMeans};
\item \texttt{KStep};
\item \texttt{FStep};
\item \texttt{CStep};
\item \texttt{KFCRegressor};
\item \texttt{KFCClassifier};
\item complete KFC training and prediction workflows.
\end{itemize}

A formal benchmarking suite should also be developed. This suite should execute the notebooks or equivalent scripts automatically, store results in a structured format, and allow fair comparison across datasets, random seeds, and model configurations.

\subsection{Improved Probability Prediction Support}

The KFC classification interface should be completed by implementing probability prediction throughout the full pipeline. This requires adding probability support to the F-step, ensuring that local models expose compatible probability outputs, and defining how probability matrices should be aggregated by the C-step.

This improvement would make the framework more useful for classification tasks where calibrated probabilities are required, such as medical risk prediction or decision-support systems.

\subsection{Improved Documentation and Reproducibility}

Future work should strengthen project documentation by adding API-level documentation, user guides, and experiment reproduction instructions. A documentation generator such as Sphinx or MkDocs could be introduced to build structured documentation from docstrings and examples.

Notebook execution could also be integrated into the validation workflow so that experimental reproducibility is checked automatically when the framework changes.

\subsection{Possible Deep Learning Integration}

Deep learning models are not part of the implemented KFCProcedure core workflow. However, future versions could investigate deep learning integration as an extension. For example, neural networks could be wrapped as local F-step estimators, or neural aggregation models could be explored as C-step combiners. This should be treated as future research only and not as an existing contribution of the current repository.

\section{Final Remarks}
\label{sec:final_remarks}

The \texttt{kfc\_procedure} framework provides a modular foundation for studying clusterwise learning and COBRA-based ensemble aggregation within a unified Python package. Its main strength lies in its flexible architecture: clustering, local learning, and aggregation are implemented as separable and configurable stages.

The repository demonstrates how divergence-based clustering, local supervised learning, and prediction-space aggregation can be connected through a scikit-learn-style workflow. At the same time, the experimental evidence remains limited to the notebooks provided, and no broad claim of consistent performance superiority is made.

Overall, this work contributes a reusable and extensible research framework for hybrid ensemble learning. With further improvements in testing, benchmarking, scalability, probability prediction, and documentation, the framework can serve as a stronger platform for future research on clusterwise modeling and adaptive ensemble aggregation.
